\documentclass{book}

\title{Parsing Foundations}
\newcommand{\booksubtitle}{A Study of Data and Parsing}
\newcommand{\booklicense}{Creative Commons Zero 1.0 Universal}

\author{Richard R. Masters}
\newcommand{\authorsubtitle}{Seattle, WA}

% Create convenient commands \booktitle and \bookauthor
\makeatletter
\newcommand{\booktitle}{\@title}
\newcommand{\bookauthor}{\@author}
\makeatother

\newcommand{\newterm}[2]{\textit{#1}\index{#2}}

% This utf8 declaration is not needed for versions of latex > 2018 but may
% be helpful for older software. Eventually it may not be worth keeping.
\usepackage[utf8]{inputenc}  
\usepackage{hyperref} % clickable links
\usepackage{fix-cm}   % this package allows large \fontsize
\usepackage{tikz}     % this is for graphics. e.g. rectangle on title page
\usepackage{amsmath}  % Used by equations
\usepackage{hyperref} % clickable links

% The following dimensions specify 4.75" X 7.5" content on 6 3/8" by 9 1/4"
% paper. The paper width and height can be tweaked as required and the content
% should size to fit within the margins accordingly.
%
% The (inside) bindingoffset should be larger for books with more pages. Some
% standard recommended sizes are .375in minimum up to 1in for 600+ page books.
% Sizes .75in and .875in are also recommended roughly at 150 and 400 pages.
\usepackage[bindingoffset=0.625in,
            left=.5in, right=.5in,
            top=.8125in, bottom=.9375in,
            paperwidth=6.375in, paperheight=9.25in]{geometry}
% Here is an alternative geometry for reading on letter size paper:
% \usepackage[margin=.75in, paperwidth=8.5in, paperheight=11in]{geometry}

\renewcommand{\contentsname}{Table of Contents} % default is {Contents}
\usepackage{makeidx}
\makeindex % Initialize an index so we can add entries with \index

% Content Starts Here
\begin{document}
\frontmatter

% ---- Half Title Page ----
% current geometry will be restored after title page
\newgeometry{top=1.75in,bottom=.5in}
\begin{titlepage}
\begin{flushleft}

% Title
\textbf{\fontfamily{qcs}\fontsize{48}{54}\selectfont Parsing\\Foundations\\}

% Draw a line 4pt high
\par\noindent\rule{\textwidth}{4pt}\\

% Subtitle
% Shaded box from left to right. Text node is midway (centered).
\begin{tikzpicture}
\shade[bottom color=lightgray,top color=white]
    (0,0) rectangle (\textwidth, 1.5)
    node[midway] {\textbf{\large \textit{\booksubtitle}}};
\end{tikzpicture}

% Edition Number
\begin{flushright}
\Large Unfinished Draft of First Edition
\end{flushright}

% \vspace{\fill}
\vspace{\fill}

\end{flushleft}
\end{titlepage}
\restoregeometry
% ---- End of Half Title Page ----


% No page numbers on the Frontispiece page
\thispagestyle{empty}


% ---- Title Page ----
% current geometry will be restored after title page
\newgeometry{top=1.75in,bottom=.5in}
\begin{titlepage}
\begin{flushleft}

% Title
\textbf{\fontfamily{qcs}\fontsize{48}{54}\selectfont Parsing\\Foundations\\}

% Draw a line 4pt high
\par\noindent\rule{\textwidth}{4pt}\\

% Shaded box from left to right with Subtitle
% The text node is midway (centered).
\begin{tikzpicture}
\shade[bottom color=lightgray,top color=white]
    (0,0) rectangle (\textwidth, 1.5)
    node[midway] {\textbf{\large \textit{\booksubtitle}}};
\end{tikzpicture}

% Edition Number
\begin{flushright}
\Large Unfinished Draft of First Edition
\end{flushright}

\vspace{\fill}

% Author and Location
\textbf{\large \bookauthor}\\[3.5pt]
\textbf{\large \textit{\authorsubtitle}}

\vspace{\fill}

% Self Publishing Logo. Free to use: CC0 license.
% The source file is book.svg. If you change the svg, you must then convert
% it to pdf. There are many online and offline tools available to do that.
\begin{center}
\includegraphics{booksvg.pdf}\\[4pt]
\fontfamily{lmtt}\small{Self Publishers Worldwide\\
Seattle San Francisco New York\\
London Paris Rome Beijing Barcelona}
\end{center}

\end{flushleft}
\end{titlepage}
\restoregeometry
% ---- End of Title Page ----

% Do not show page numbers on colophon page
\thispagestyle{empty}

\begin{flushleft}
\vspace*{\fill}
This book was typeset using \LaTeX{} software.\\
\vspace{\fill}
Copyright \textcopyright{} \the\year{}  \bookauthor\\
License: \booklicense
\end{flushleft}

% A title page resets the page # to 1, but the second title page
% was actually page 3. So add two to page counter.
\addtocounter{page}{2}

% The asterisk excludes chapter from the table of contents.
\chapter*{Preface}
% 0000000100000000020000000003000000000400000000050000000006000000000700000000008
This book was created from a desire to make parsing easier to understand.
Parsing is an important topic in the field of computer science and the author
struggled to learn about parsing. It was not that the ideas themselves were
particularly difficult, but that it required a huge investment of time and some
money to identify, locate, acquire, organize, and understand the relevant
information. In many cases, the information presented assumes the reader
is fluent with set-builder notation\footnote{\url{https://en.wikipedia.org/wiki/Set-builder_notation}}
and sometimes particular symbols may not be defined because they have been
used in previous papers on the subject that the reader is assumed to
have read.

It became clear that it would be a lot easier to learn the subject from a
book that covers the relevant topics with through explanations. It was
also clear that there is a lot of existing research, some
of it several decades old, that could benefit from a wider audience.
Finally, a good understanding of parsing rests on a good
understanding of data, which is what a parser acts upon.

\section*{Prior Work}
Prior to the writing of this book, there were not that many resources for
learning about parsing that were both freely available and of high-quality.
There were no more than a handful of books about parsing on the market and
the most relevant to a general scientific audience was a relatively expensive
book about 
compilers\footnote{\url{https://en.wikipedia.org/wiki/Compilers:_Principles,_Techniques,_and_Tools}}.
There are also quite a few papers and articles written for academic
journals that are written for a narrow audience by use of vocabulary and notation.
Finally, while there are some good online resources, they are lost within a sea
of information that is either outdated or poorly produced, or only covers a
limited set of topics. 

\section*{Design Goals}
% 0000000100000000020000000003000000000400000000050000000006000000000700000000008
This book is designed to be as simple to understand as possible. 
This book is built from the ground up with relevant topics,
starting from the basics and building up to advanced topics, 
using a simple and intuitive style of presentation.

Some of the explanations in this book may seem excessively long so it is
worth noting that keeping explanations short is a much lower priority
for this book than providing a thorough explanation. The idea is that far
more time is wasted by avoiding difficult material, skimming it repeatedly,
and finally focusing your mind enough to understand it, than is wasted
reading and understanding material that is excessively explained.


\section*{Audience}
This book can serve as an introduction to topics on data and parsing for
the working software developer or student of computer science. The book
covers foundational topics of parsing using concrete examples and common
% 0000000100000000020000000003000000000400000000050000000006000000000700000000008
tools used in the field. In an academic setting, the book would serve well as
a companion book for courses that only cover parsing as part of a larger
topic, such as compiling.

\section*{Uses for this Book}
The book may also be of interest to you if you are an educator, a blogger,
or if you are writing a book. This document is licensed under the Creative
Commons Zero v1.0 Universal 
license\footnote{\url{https://creativecommons.org/publicdomain/zero/1.0/legalcode}}.
You are free to copy this content,
distribute it, include it in your own content, or sell it, and you may do so
\textit{without attribution}. Just copy what you want and move on.

\section*{Terms and Symbols}
Many of the ideas covered in this book originated in the 1950s and 
1960s, near the dawn of the computer age. In the decades since,
researchers and educators have defined quite a few names for various
concepts surrounding data and parsing. They have also adopted numerous
symbols and notations as a convenient shorthand to write down their ideas.
% 0000000100000000020000000003000000000400000000050000000006000000000700000000008
While it can be tedious to learn (and explain) these names and symbols, it
is necessary if you want to be able to follow discussions or read
original research papers on these topics. In some cases, words that you may
be familiar with in common language may have a different definition in 
academic papers. This can be confusing if you are not already familiar with
that term. Therefore, in this book, we will error on the side of explaining
the academic terms and symbols when they are used. Sometimes it takes a 
number of definitions conspiring together to explain a larger concept, so
bear with us if it seems like we are just defining words for no apparent
reason.

\section*{Prerequisites}
For the chapters on data, the reader should have a basic knowledge of
computers, such as knowing about a monitor, mouse, disk, file, and an
application.

% Three-level Table of Contents
\setcounter{tocdepth}{3}
\tableofcontents

\mainmatter

\chapter{Data}\label{ch.Data}

\section{Symbols}
% 0000000100000000020000000003000000000400000000050000000006000000000700000000008
The signs, sounds, marks, and words humans use to represent ideas are called
\newterm{symbols}{Symbol}. For example, the twenty six letters of
the English alphabet are symbols representing sounds or parts of words. Words
are also symbols for the ideas they represent. 
The numerals 0 through 9 are symbols for numbers.
As long as something represents, or \textit{symbolizes} something else, 
it is a symbol. We use symbols as sort of "real world" currency
to remember or communicate something meaningful to ourselves or others.

\subsection{Interpreting Symbols}
Wikipedia defines
\newterm{data}{Data} as any sequence of symbols for which "a specific act of
interpretation" gives those symbols
meaning\footnote{\url{https://en.wikipedia.org/wiki/Data\_(computing)}}.
In other words, the essential nature of data is that is in a
\newterm{sequence}{Sequence}, which means that the symbols occur one after
another, and that there is a particular meaning associated with 
the ordering and choice of the symbols in the data.

\subsection{Manipulating Symbols}

% 0000000100000000020000000003000000000400000000050000000006000000000700000000008
Computers are well known for their ability to manipulate data.
A computer can read and \newterm{execute}{Execution}, or follow,
instructions as represented by a sequence of symbols.
These instructions can, in
turn, direct the manipulation of symbols stored in the computer's memory.
It can be confusing to think of a computer working with the kind of visual
symbols that humans typically think of like a + (plus) sign or the numeral 2,
and that would be because they do not actually work directly with those
kinds of symbols.
The symbols that a computer can read, follow, and manipulate are quite limited.
In fact at the lowest levels, there are only two symbols. Inside the computer,
these symbols exist as electrical states inside the transistors and wires of 
the computer. From an electrical perspective, we think of these states as 
\textit{off} and \textit{on}, but from a numerical perspective these states 
symbolize 0 and 1, respectively.

\section{Binary Data}\label{bindata}
\subsection{Bits}
A numerical symbol which is either 0 or 1 is known as a
\newterm{binary digit}{Binary Digit}, or \newterm{bit}{Bit} for short.
At the most fundamental level, bits are the letters of a very simple
alphabet that both humans and computers can manipulate in an
equivalent way. For example, both humans and computers are capable 
of numerically adding together the symbols 1 and 0 to produce the 
symbol 1. It is through these common definitions and ways of processing
symbols that computers can perform computations
that are relevant to humans.

% 0000000100000000020000000003000000000400000000050000000006000000000700000000008
\subsection{Binary Codes}
Since computers are only able to manipulate sequences of bits, we
have devised an incredibly rich variety of ways, known as 
\newterm{codes}{Codes}\footnote{\url{https://en.wikipedia.org/wiki/Code}}
to map the information available to us into sequences of ones and zeros. For
example, \newterm{decimal}{Decimal}\footnote{\url{https://en.wikipedia.org/wiki/Decimal}} 
numbers, which are the standard numbers we use from day to day, can be converted
to a binary number composed of binary digits, using a
\newterm{binary code}{Binary Code} as seen in Table 1.1.


\begin{table}[h!]
\centering
\begin{tabular}{||c|c||}
    \hline
    Number & Bits\\
    \hline\hline
     0 & 000 \\
     \hline
     1 & 001 \\
     \hline
     2 & 010 \\
     \hline
     3 & 011 \\
     \hline
     4 & 100 \\
     \hline
     5 & 101 \\
     \hline
     6 & 110 \\
     \hline
     7 & 111 \\
     \hline
\end{tabular}
\caption{Mapping of Numbers to Bits}
\end{table}

\subsection{Nibbles}
A \newterm{nibble}{Nibble} is four bits.

\subsection{Hexadecimal}
Hexadecimal is a numerical encoding that allows the use of
16 different symbols for each digit, 0 through 9 and A through F.
Each symbol encodes 4 bits of a binary number.

\subsection{Bytes and Octets}
A \newterm{byte}{Byte} is typically 8 bits. An \newterm{octet}{Octet} is 
always 8 bits.

% 0000000100000000020000000003000000000400000000050000000006000000000700000000008
\subsection{Exercises for Section \ref{bindata}}\label{ex.bindata}

\noindent\textbf{Exercise \ref{ex.bindata}.1}: Write the numbers 8 through 15
using 4 bits.

\medskip\noindent
\textbf{Exercise \ref{ex.bindata}.2}: Create a mapping from the numbers 0
through 8 to two digits of the three symbols 0, 1, and 2.

\section{Character Encodings}
Characters are symbols that represent sounds or parts of words.

\subsection{ASCII}

\subsection{Terminals and Control Characters}

\subsection{Locales}

\subsection{Unicode}
The Unicode\textcircled{r} Standard is a universal character encoding standard for all written characters and text.

\section{Summary of Chapter \ref{ch.Data}}

\chapter{Parsing}

% 0000000100000000020000000003000000000400000000050000000006000000000700000000008
\newterm{Parsing}{Parsing} as a way of interpreting data based on how it matches
\newterm{structural patterns}{Structural Patterns}, such as a
\newterm{formal grammar}{Formal Grammar}\index{Grammar}. 

\backmatter
\addcontentsline{toc}{chapter}{Index}
\printindex
\end{document}
